% !PRINTED 680
\weberpage{680}
%<<CONT
die Exponenten $\xi_i$ auf nicht negative echte gebrochene Werthe beschr\"ankt bleiben, und daraus ergiebt sich nach der Definition der conjugirten Logarithmen eine obere Grenze f\"ur $\varepsilon$ und seine Conjugirten.

Nach dem Satze 2., \S.\,185 giebt es aber in $\Omega$ nur eine endliche Anzahl ganzer Zahlen, also um so mehr nur eine endliche Anzahl von Einheiten, die dem absoluten Werthe nach mit allen ihren Conjugirten unter einer endlichen Grenze liegen.

Wenn wir nun eine Einheit, deren Exponenten in den Grenzen 0 und 1 liegen, mit Einschluss der unteren und mit Ausschluss der oberen Grenze eine \gesperrt{in Bezug auf das System $\varepsilon_1, \varepsilon_2, \ldots, \varepsilon_{\nu-1}$ reducirte Einheit nennen}, so haben wir den Satz:

\gesperrt{Es giebt nur eine endliche Anzahl von Einheiten in $\Omega$, die in Bezug auf ein unabh\"angiges System $\varepsilon_1, \varepsilon_2, \ldots, \varepsilon_{\nu-1}$ reducirt sind.}

2) Die Einheiten reproduciren sich durch Multiplication und Division. Sind
\[
\begin{array}{l}
\xi_1', \xi_2', \ldots, \xi_{\nu-1}' \\
\xi_1'', \xi_2'', \ldots, \xi_{\nu-1}''
\end{array}
\]
die Exponenten von zwei Einheiten $\varepsilon', \varepsilon''$, so sind die Summen und die Differenzen
\[
\xi_1' \pm \xi_1'', \xi_2' \pm \xi_2'', \ldots, \xi_{\nu-1}' \pm \xi_{\nu-1}''
\]
die Exponenten der Einheiten $\varepsilon' \varepsilon''$ und $\varepsilon' : \varepsilon''$.

Dies ergiebt sich unmittelbar aus dem Satze, dass der Logarithmus eines Productes oder eines Quotienten gleich der Summe oder der Differenz der Logarithmen der beiden Bestandtheile ist.

3) Die Einheiten $\varepsilon_1, \varepsilon_2, \ldots$ selbst haben die Exponenten $1, 0, \ldots, 0; 0, 1, \ldots, 0; \ldots$, und daraus folgt nach 2), dass jedes System von ganzen Zahlen, f\"ur $\xi_1, \xi_2, \ldots, \xi_{\nu-1}$ gesetzt, das Exponentensystem einer Einheit giebt, die sich durch Multiplication von Potenzen der $\varepsilon_1, \varepsilon_2, \ldots$ bilden l\"asst.

4) Aus 2) und 3) ergiebt sich, dass ein Exponentensystem $\xi_1, \xi_2, \ldots, \xi_{\nu-1}$ einer Einheit ein Exponentensystem bleibt, wenn alle seine Elemente mit einer und derselben ganzen Zahl multiplicirt werden, und dass es auch dann noch diesen Charakter beh\"alt, wenn seine Elemente $\xi_i$ um beliebige ganze Zahlen vermehrt oder vermindert werden. Gebrauchen wir also das Zeichen $(x)$ in dem Sinne, wie im \S.\,182, dass es den Ueber-
%CONT>>
